\documentclass{tum_report}
\usepackage[T1]{fontenc}
\usepackage{combelow}
\usepackage{booktabs}
\usepackage{tabularx}
\usepackage{microtype}

\hypersetup{pdftitle={Laboratory Work No. 1: Caesar Cipher},
    pdfauthor={Mihail Pancenco}, pdfsubject={Cryptography, Romanian alphabet}}
\lstdefinestyle{cryptocode}{
    language=[Sharp]C, basicstyle=\ttfamily\fontsize{9}{11}\selectfont,
    keywordstyle=\bfseries\color{blue!45!black},
    commentstyle=\itshape\color{black!55}, stringstyle=\color{green!35!black},
    numbers=left, numberstyle=\tiny\color{black!55}, numbersep=7pt,
    frame=single, rulecolor=\color{black!25}, backgroundcolor=\color{black!2},
    breaklines=true, breakatwhitespace=false, columns=fullflexible,
    keepspaces=true, showstringspaces=false, tabsize=4, emptylines=0,
    aboveskip=6pt, belowskip=6pt, xleftmargin=12pt, xrightmargin=0pt,
    literate={Ă}{{\u{A}}}1 {ă}{{\u{a}}}1 {Â}{{\^{A}}}1 {â}{{\^{a}}}1
      {Î}{{\^{I}}}1 {î}{{\^{i}}}1 {Ș}{{\cb{S}}}1 {ș}{{\cb{s}}}1
      {Ț}{{\cb{T}}}1 {ț}{{\cb{t}}}1 {Ş}{{\c{S}}}1 {ş}{{\c{s}}}1
      {Ţ}{{\c{T}}}1 {ţ}{{\c{t}}}1
}
\newcommand{\codepart}[3]{{\setstretch{1}\lstinputlisting[style=cryptocode,
    firstline=#1,lastline=#2,gobble=#3]{../code/Program.cs}}}
\newcommand{\screenshot}[3]{%
    \begin{figure}[H]\centering
    \includegraphics[width=#3\textwidth]{img/#1.png}
    \caption{#2}\label{fig:#1}\end{figure}}

\begin{document}
\ministry{MINISTRY OF EDUCATION AND RESEARCH\\OF THE REPUBLIC OF MOLDOVA}
\university{\textbf{Technical University of Moldova}}
\faculty{\textbf{Faculty of Computers, Informatics and Microelectronics}}
\department{\textbf{Department of Software Engineering and Automatics}}
\title{Cryptography}{}
\year{2026}

\begin{titlepage}
\centering
\includegraphics[width=0.25\textwidth]{UTM_Logo.png}\par\vspace{0.8cm}
{\bfseries
\textsc{\ministryname}\\[0.15cm]
\universityname\\
\facultyname\\
\departmentname\par}

\vspace{2cm}
\textbf{Mihail Pancenco}\quad\textbf{FAF-243}\par
\vspace{0.9cm}
{\Large\textbf{Report}\par}
\vspace{0.5cm}
{\Large\normalfont\textit{Laboratory Work No. 1}\par}
\vspace{0.35cm}
{\Large\normalfont\textit{Cryptography}\par}
\vspace{0.7cm}
{\huge\normalfont\textit{Caesar Cipher}\par}
\vspace{0.4cm}
{\large\normalfont Standard and two-key versions\\over the Romanian alphabet\par}

\vspace{2cm}
\begin{flushright}
\normalfont
Supervisors:\par
\textbf{Poștaru Andrei}, \textit{University lecturer}\par
\textbf{Maia Zaica}, \textit{University assistant}\par
DISA, FCIM, UTM
\end{flushright}
\vfill
\textbf{Chișinău -- \degreeyear}\par
\end{titlepage}

\clearpage
\addtocontents{toc}{\protect\thispagestyle{empty}}
\tableofcontents
\thispagestyle{empty}
\clearpage
\pagestyle{fancy}

\chapter{Short theory}
\section{Objectives and the Romanian alphabet}
The aim of Laboratory Work No. 1 is to implement the standard Caesar cipher
and a Caesar cipher with a keyword-based alphabet permutation, to understand
modular encryption and decryption, and to compare their security. Both tasks
use the 31-letter Romanian alphabet and the encoding specified in the
laboratory assignment~\cite{ref:assignment}.

Each letter is represented by its position, starting at zero. Table~\ref{tab:alphabet}
shows the complete encoding. Character codes such as ASCII or Unicode values
are not used to calculate the shift.
\begin{table}[H]
\caption{Romanian alphabet and the required numeric encoding}
\label{tab:alphabet}
\centering\small
\setlength{\tabcolsep}{5pt}
\begin{tabular}{l*{16}{c}}\toprule
Letter & A & Ă & Â & B & C & D & E & F & G & H & I & Î & J & K & L & M\\
Code & 0 & 1 & 2 & 3 & 4 & 5 & 6 & 7 & 8 & 9 & 10 & 11 & 12 & 13 & 14 & 15\\\midrule
Letter & N & O & P & Q & R & S & Ș & T & Ț & U & V & W & X & Y & Z & \\
Code & 16 & 17 & 18 & 19 & 20 & 21 & 22 & 23 & 24 & 25 & 26 & 27 & 28 & 29 & 30 & \\\bottomrule
\end{tabular}
\end{table}

\section{Task 1.1: the standard Caesar cipher}
A Caesar cipher replaces every plaintext letter with the letter located a
fixed number of positions further along the alphabet. The secret key is an
integer $k\in\{1,2,\ldots,30\}$. For a plaintext index $x$ and a ciphertext
index $y$, encryption and decryption are
\begin{align}
 y &= E_k(x)=(x+k)\bmod 31,\\
 x &= D_k(y)=(y-k+31)\bmod 31.
\end{align}
Adding 31 during decryption prevents a negative remainder. This is sufficient
because the accepted key is at most 30. For example, with $k=3$, the letter
C has index 4 and becomes F, at index 7; Z has index 30 and becomes Â,
at index 2. Decrypting  gives $(2-3+31)\bmod31=30$, which is Z.

The standard cipher has exactly 30 permitted keys; $k=0$ is excluded because
it would leave the text unchanged. An attacker can decrypt the ciphertext
with each of the 30 keys and inspect the candidates. This small key space
makes exhaustive search practical.

\clearpage
\section{Task 1.2: Caesar cipher with two keys}
The second version uses a numeric key $k_1\in\{1,\ldots,30\}$ and a keyword
$k_2$ containing at least seven Romanian letters. The keyword is converted
to uppercase. Its letters are placed first, retaining only the first
occurrence of each letter, and the unused letters of the Romanian alphabet
are appended in their original order. The resulting alphabet $P$ contains
all 31 letters exactly once.

If $\operatorname{pos}_P(a)$ denotes the position of letter $a$ in $P$, the
letter-level formulas are
\begin{align}
 E_{k_1,P}(a)&=P[(\operatorname{pos}_P(a)+k_1)\bmod31],\\
 D_{k_1,P}(b)&=P[(\operatorname{pos}_P(b)-k_1+31)\bmod31].
\end{align}
Both the lookup of the input letter and the lookup of the output letter use
the same permuted alphabet. The keyword changes the alphabet order, while
the numeric key determines the cyclic shift.

For $k_2=\texttt{CRIPTOGRAFIE}$, the distinct keyword letters are
\texttt{CRIPTOGAFE}. Appending the remaining letters gives
\begin{center}\small\ttfamily
CRIPTOGAFEĂÂBDHÎJKLMNQSȘȚUVWXYZ
\end{center}
With $k_1=3$, C at position 0 becomes P at position 3, while I at position 2
becomes O at position 5.

Considering every possible alphabet permutation gives the theoretical
number of parameter pairs
\[
 |K_{\mathrm{pairs}}|=30\cdot31!\approx2.47\times10^{35}.
\]
This counts pairs of a shift and an alphabet ordering, rather than keyword
strings: repeated letters can make different keywords produce the same
ordering. Different parameter pairs can also produce the same substitution,
so the number of distinct encryptions is at most $31!$. Restricting keywords
to ordinary words reduces the practical search space further.

Despite this larger theoretical key space, the cipher remains
monoalphabetic: a given plaintext letter always maps to the same ciphertext
letter. Letter counts and repeated-letter patterns are preserved under
renaming, so frequency analysis can still be used. Neither version provides
security suitable for protecting modern confidential information.

\clearpage
\chapter{Implementation}
\section{Program structure and letter encoding}
The implementation is a C\# console program in \texttt{lab\_01/code/Program.cs}.
\texttt{RomanianAlphabet} provides shared normalization, validation, letter
conversion and modular shifting. \texttt{CaesarCipher} implements Task 1.1,
\texttt{TwoKeyCaesarCipher} implements Task 1.2, and
\texttt{CipherTablePrinter} displays the numeric and letter mappings.
\texttt{Program} manages the menu and input retry loops. The following
excerpts are read directly from the existing source file; empty lines are
compressed for readability.

\codepart{10}{11}{4}
\noindent\textit{Comment.} The alphabet string exactly matches Table~\ref{tab:alphabet},
and its size is explicitly set to 31.

\codepart{75}{85}{8}
\noindent\textit{Comment.} \texttt{IndexOf} converts each letter to its position
in the supplied alphabet. For Task 1.1 this is the original alphabet; for
Task 1.2 it is the permuted alphabet. A missing letter causes an exception,
rather than an invalid numeric index.

\codepart{90}{112}{4}
\noindent\textit{Comment.} \texttt{ToText} maps the shifted indices back to
letters. \texttt{ShiftNumbers} implements modular arithmetic. Encryption
passes a positive key; decryption passes its negative. Adding \texttt{Size}
keeps the remainder nonnegative for the validated key range.

\clearpage
\section{Normalization and message validation}
\codepart{13}{22}{4}
\noindent\textit{Comment.} Normalization replaces the cedilla variants
\texttt{ş/Ş} and \texttt{ţ/Ţ} with the Romanian comma-below letters
\texttt{ș/Ș} and \texttt{ț/Ț}, then converts the string to uppercase.
This supports both uppercase and lowercase input from different keyboards.

\codepart{24}{40}{4}
\noindent\textit{Comment.} Validation accepts only letters from the explicit
alphabet. Ordinary spaces are allowed in messages. An error identifies the
rejected character and lists the allowed letters. Digits, punctuation,
tabs and other characters are rejected.

\noindent The helper \texttt{Prepare} normalizes and validates the input
before removing all ordinary spaces. Both encryption and decryption use
this preparation. Consequently, decryption recovers uppercase text without
spaces; the original capitalization and word spacing cannot be reconstructed.

\codepart{50}{61}{4}
\noindent\textit{Comment.} Keywords are normalized in the same way as messages,
but spaces are forbidden and the normalized keyword must contain at least
seven letters. This length check occurs before duplicate removal, so the
keyword does not need seven \emph{distinct} letters.

\clearpage
\section{Building and displaying the permuted alphabet}
\codepart{185}{210}{4}
\noindent\textit{Comment.} The first loop keeps the first occurrence of each
keyword letter. The second loop appends every unused Romanian letter in
the prescribed order. As a result, the alphabet remains a permutation of
all 31 letters even when the keyword contains duplicates or diacritics.

\codepart{212}{223}{4}
\noindent\textit{Comment.} Task 1.2 displays the normalized keyword, the original
alphabet and the permuted alphabet, making the ordering available for manual
verification. Both cipher classes also print an encryption table with four
rows: \texttt{x}, \texttt{m}, \texttt{x+k} and \texttt{c}. The table always
shows the forward mapping, including when the selected operation is decryption;
the inverse operation itself uses the negative shift.

\clearpage
\section{Encryption and decryption}
\codepart{121}{143}{4}
\noindent\textit{Comment.} The standard cipher first validates the numeric key
and prepares the message. It then converts letters to positions in the
original alphabet, applies $+k$ for encryption or $-k$ for decryption, and
converts the resulting positions back to letters.

\codepart{159}{183}{4}
\noindent\textit{Comment.} The two-key cipher follows the same steps, but builds
the keyword alphabet first. Both \texttt{ToNumbers} and \texttt{ToText}
receive that alphabet. Decryption must use the same numeric key and the same
keyword-derived ordering as encryption.

\clearpage
\section{Numeric keys, retry loops and the console menu}
\codepart{63}{69}{4}
\noindent\textit{Comment.} The cipher methods independently enforce the
numeric key interval $1\leq k\leq30$, so invalid keys are rejected even if
a method is called outside the console menu.

\codepart{348}{363}{4}
\noindent\textit{Comment.} The console accepts a key only if
\texttt{int.TryParse} succeeds and the value lies between 1 and 30.
Otherwise it displays the permitted range and asks for the key again.

\codepart{365}{382}{4}
\noindent\textit{Comment.} The keyword reader catches validation exceptions,
displays the reason and repeats the prompt. \texttt{ReadMessage} uses the
same pattern with \texttt{Prepare}, so invalid message characters also cause
a retry. \texttt{ReadChoice} validates the available menu options.

The menu selects a cipher, then encryption or decryption. The user enters
the numeric key, the keyword when required, and the message. The program
prints the mapping table, prepared input and result, then returns to the
menu; option 0 exits. The substitution uses explicit alphabet indices.

\clearpage
\chapter{Results}
\section{Task 1.1: encryption and decryption}
The unchanged program was compiled with the .NET 10 SDK. The screenshots
show a live local console viewer connected to the program's standard input
and output. The viewer only displays the console session; all cipher logic
and validation come from \texttt{Program.cs}.

For the input \texttt{cifrul cezar} and $k=3$, the prepared plaintext is
\texttt{CIFRULCEZAR}. Encryption produces \texttt{FKITXOFHÂBT}; decrypting
this ciphertext with the same key returns \texttt{CIFRULCEZAR}. These
outputs use the Romanian alphabet, so they differ from the English-alphabet
example in the assignment.
\screenshot{standard_encrypt}{Task 1.1: encryption with the numeric key 3}{1.0}
\noindent\textit{Manual check.} The first three letters of \texttt{CIFRULCEZAR}
are C, I and F, with positions 4, 10 and 7. Adding 3 gives positions 7, 13
and 10, or F, K and I. The Z at position 30 becomes  at position 2.

\clearpage
\subsection{Recovering the standard-cipher plaintext}
The user selects the standard cipher and operation 2, then supplies key 3
and ciphertext \texttt{FKITXOFHÂBT}. The result in
Figure~\ref{fig:standard_decrypt} is \texttt{CIFRULCEZAR}, matching the
prepared input from the encryption run.
\screenshot{standard_decrypt}{Task 1.1: decryption with the numeric key 3}{1.0}
\noindent\textit{Manual check.} F at position 7 becomes C at position
$(7-3+31)\bmod31=4$. K at position 13 becomes I at position 10.
 at position 2 becomes Z at position $(2-3+31)\bmod31=30$.
Spaces and lowercase letters are absent because preparation removes them.

\clearpage
\section{Task 1.2: encryption and decryption}
For $k_1=3$ and the lowercase keyword \texttt{criptografie}, the program
displays \texttt{CRIPTOGRAFIE} and the alphabet
\texttt{CRIPTOGAFEĂÂBDHÎJKLMNQSȘȚUVWXYZ}. Encrypting
\texttt{cifrul cezar} produces \texttt{POÂTXQPBIĂT}. Decryption with the
same two keys restores \texttt{CIFRULCEZAR}.

The printed table also permits manual checking: C maps to P, I to O,
F to Â, and Z wraps from position 30 to position 2, which contains I.
\screenshot{two_key_encrypt}{Task 1.2: encryption and the displayed permuted alphabet}{1.0}
\noindent\textit{Comparison.} With the same numeric key and prepared input,
the standard version gives \texttt{FKITXOFHÂBT}, while the permuted version
gives \texttt{POÂTXQPBIĂT}. The difference comes entirely from the
keyword-derived alphabet order.

\clearpage
\subsection{Recovering the two-key plaintext}
To decrypt, the user selects mode 2 and operation 2. The numeric key is 3
and the keyword is \texttt{criptografie}. The ciphertext is
\texttt{POÂTXQPBIĂT}. The program rebuilds the same displayed permutation
and returns \texttt{CIFRULCEZAR}.
\screenshot{two_key_decrypt}{Task 1.2: decryption using the same numeric key and keyword}{1.0}
\noindent\textit{Manual check.} In the permuted alphabet, P has position 3
and decrypts to position 0, containing C. O has position 5 and decrypts to
position 2, containing I. I in the ciphertext has position 2 and decrypts
to position 30, containing Z. These checks confirm that decryption also
uses positions in the permuted alphabet.

\clearpage
\section{Invalid numeric keys and message characters}
Figure~\ref{fig:invalid_input} shows three rejected numeric keys:
\texttt{0}, \texttt{31} and \texttt{abc}. Each attempt produces
\texttt{Invalid key. Enter an integer between 1 and 30.}, followed by a
new key prompt. Once the valid key \texttt{3} is supplied, the program
continues to the message.

The messages \texttt{salut!} and \texttt{salut2} are rejected. Each error
names the offending character, displays the accepted alphabet and states
that spaces are allowed. The message prompt is repeated after each error.
A subsequent valid input \texttt{Salut} was accepted and encrypted to
\texttt{ȚBOXV}, confirming that the program recovers from invalid input.
\screenshot{invalid_input}{Rejected keys and message characters, followed by repeated prompts}{1.0}

\clearpage
\section{Invalid keywords}
Figure~\ref{fig:invalid_keyword} shows that the two-key mode also rejects the
numeric key \texttt{-1}. The key \texttt{30} is accepted. The six-letter
keyword \texttt{secret} produces the minimum-length error;
\texttt{cripto1} is rejected because it contains a digit, and
\texttt{cripto grafie} is rejected because keywords cannot contain spaces.
Each error is followed by a new keyword prompt.
\screenshot{invalid_keyword}{Two-key validation: wrong shift, short keyword, digit and space}{1.0}

\clearpage
\section{Romanian letters and boundary checks}
With the standard cipher, $k=1$ and input \texttt{aăâîşţz}, normalization
produces \texttt{AĂÂÎȘȚZ}; encryption produces \texttt{ĂÂBJTUA}.
Figure~\ref{fig:romanian_letters} demonstrates lowercase conversion,
cedilla compatibility, Romanian diacritics and the wraparound $Z\to A$.
\screenshot{romanian_letters}{Romanian letters, cedilla normalization and wraparound at key 1}{0.95}

The two-key mode was also checked with $k_1=30$, the seven-letter keyword
\texttt{ştiinţă} and the same input. The normalized keyword
\texttt{ȘTIINȚĂ} yields \texttt{ȘTINȚĂAÂBCDEFGHÎJKLMOPQRSUVWXYZ};
the result is \texttt{ĂȚAHZNY}. This checks the upper key boundary and
shows that a valid seven-letter keyword can include a repeated letter.

Additional checks exercised every numeric key from 1 to 30 over all 31
letters: once for the standard cipher and once for each of the keywords
\texttt{criptografie} and \texttt{ştiinţă}. All 90 encryption/decryption
round trips passed. Each encryption was also compared with an independently
calculated modular mapping. In total, these checks executed the unchanged
program 180 times. Console transcripts and the verification summary are
saved in the report's \texttt{evidence} folder.

\clearpage
\chapter{Conclusions}
In this laboratory work, I implemented and examined two Caesar ciphers for
the Romanian alphabet: a fixed-alphabet version and a version with a
keyword-based permutation. I used explicit positions from 0 to 30, rather
than character encoding values, and applied modular arithmetic to encrypt
and decrypt the text. The results confirmed that decryption with the same
key or keys recovers the prepared plaintext.

The shared normalization and validation functions handle Romanian
diacritics, cedilla variants, lowercase letters and spaces in messages.
The console retry loops reject invalid keys, unsupported characters and
keywords that are too short. The displayed permutation and mapping tables
make both algorithms easy to verify by hand. These checks helped me
understand why alphabet order and consistent preprocessing are essential
for correct encryption and decryption.

The standard version is weak because an attacker needs to try only 30
numeric keys. A keyword permutation greatly increases the theoretical
parameter space, but the second version still uses one fixed substitution
for the entire message and preserves letter frequencies and repeated
patterns. Therefore, neither cipher is suitable for modern secure
communication. The main value of these implementations is learning
substitution ciphers, modular arithmetic and reliable input handling.

\section{Source code}
The full source code is available in the GitHub repository~\cite{ref:source}:
\begin{center}
\url{https://github.com/MiHCraft/Cryptography_UTM}
\end{center}
The implementation for this laboratory work is
\texttt{lab\_01/code/Program.cs}. The report describes and quotes this file;
no changes were made to the implementation while preparing the report.

\clearpage
\phantomsection
\addcontentsline{toc}{chapter}{Bibliography}
\begin{thebibliography}{9}
\singlespace\normalsize
\bibitem{ref:assignment}
Technical University of Moldova.
\textit{Laboratory Work No. 1: Caesar Cipher}.
Course assignment supplied as
\texttt{Laboratory Work no. 1. Caesar Cipher.pdf}, pp.~1--4.

\bibitem{ref:source}
Pancenco, Mihail.
\textit{Cryptography\_UTM}. Romanian Caesar Cipher implementation.
\texttt{lab\_01/code/Program.cs}.
\url{https://github.com/MiHCraft/Cryptography_UTM}.
\end{thebibliography}

\end{document}
